
\documentclass{article}
\usepackage{colortbl}
\usepackage{makecell}
\usepackage{multirow}
\usepackage{supertabular}

\begin{document}

\newcounter{utterance}

\centering \large Interaction Transcript for game `dond', experiment `coop\_en', episode 6 with reflexion{-}Qwen3.5{-}4B.
\vspace{24pt}

{ \footnotesize  \setcounter{utterance}{1}
\setlength{\tabcolsep}{0pt}
\begin{supertabular}{c@{$\;$}|p{.15\linewidth}@{}p{.15\linewidth}p{.15\linewidth}p{.15\linewidth}p{.15\linewidth}p{.15\linewidth}}
    \# & $\;$A & \multicolumn{4}{c}{Game Master} & $\;\:$B\\
    \hline

    \theutterance \stepcounter{utterance}  
    & & \multicolumn{4}{p{0.6\linewidth}}{
        \cellcolor[rgb]{0.9,0.9,0.9}{
            \makecell[{{p{\linewidth}}}]{
                \texttt{\tiny{[P1$\langle$GM]}}
                \texttt{You are playing a cooperative negotiation game in which you have to agree on how to divide a set of items among you and another player.} \\
\\ 
\texttt{Rules:} \\
\texttt{(a) You and the other player are given a set of items. Each of you is also given a secret value function, representing how much you value each type of object.} \\
\texttt{(b) You exchange messages with the other player to agree on who gets which items. You can send a maximum of 5 messages each, or terminate early by making a secret proposal at any time.} \\
\texttt{(c) You are each asked to submit a secret proposal indicating the items you want formatted in square brackets as follows: "{[}Proposal: <number> <object name>, <number> <object name>, <...>{]}"} \\
\texttt{(d) If your proposals are complementary, i.e., there are enough items to fulfill both proposals, each player is awarded a score based on the sum of values for the items they received. Otherwise, both of you get zero points.} \\
\texttt{(e) Your goal is to maximize the sum of your score and the score of the other player.} \\
\\ 
\texttt{Let us start.} \\
\\ 
\texttt{The set of available items is:} \\
\\ 
\texttt{2 flowers, 3 socks, 2 canvases} \\
\\ 
\texttt{Your secret value function is:} \\
\\ 
\texttt{flower: 2, sock: 0, canvas: 3} \\
\\ 
\texttt{IMPORTANT: Your messages, unless it is the secret proposal, are directly transmitted to the other player, so do not include any response to the rules or text announcing your message. To make a secret proposal, use the indicated format. Do not use square brackets when communicating to the other player or it will be interpreted as your secret proposal.} \\
            }
        }
    }
    & \\ \\

    \theutterance \stepcounter{utterance}  
    & \multicolumn{4}{p{0.6\linewidth}}{
        \cellcolor[rgb]{0.9,0.9,0.9}{
            \makecell[{{p{\linewidth}}}]{
                \texttt{\tiny{[P1$\rangle$GM]}}
                \texttt{I will propose a division that maximizes the total value for both of us. Since I value canvases the most (3 points each) and flowers second (2 points each), and socks are worthless to me (0 points), I should aim to get the canvases and flowers for myself, while giving the socks to the other player. This allows the other player to get the socks (which they might value) while I get the high{-}value items.} \\
\\ 
\texttt{Let's assume the other player values socks highly (since they are worthless to me, they are likely the only ones who want them).} \\
\texttt{My plan:} \\
\texttt{{-} I take: 2 flowers, 2 canvases.} \\
\texttt{{-} Other player takes: 3 socks.} \\
\\ 
\texttt{Total value for me: 2*2 + 2*3 = 4 + 6 = 10.} \\
\texttt{Total value for other player: 3 * (value of sock). Even if they value socks at 1, total is 13. If they value them higher, total is even higher.} \\
\texttt{Alternative: I take 2 flowers, 1 canvas. Other takes 3 socks, 1 canvas.} \\
\texttt{My value: 4 + 3 = 7.} \\
\texttt{Other value: 3*sock + 3.} \\
\texttt{Total: 10 + 3 + 3 = 16 (assuming sock=1).} \\
\texttt{My previous plan (I take 2 canvases) gives me 10. If other takes 3 socks, and sock=1, they get 3. Total 13.} \\
\texttt{Wait, if I take 2 canvases and 2 flowers, that's 4 items. The other player gets 3 socks.} \\
\texttt{Is there a better split?} \\
\texttt{If I take 2 canvases (6 pts) and 2 flowers (4 pts) = 10 pts.} \\
\texttt{Other gets 3 socks.} \\
\texttt{If the other player values socks at 1 each, they get 3. Total 13.} \\
\texttt{If I take 1 canvas (3 pts) and 2 flowers (4 pts) = 7 pts.} \\
\texttt{Other gets 2 canvases (6 pts) and 3 socks (3 pts) = 9 pts. Total 16.} \\
\texttt{This seems better if the other player values canvases too. But I don't know their values.} \\
\texttt{However, the rule says "maximize the sum of your score and the score of the other player".} \\
\texttt{If I give the other player the canvases, they might value them highly. If I keep them, they might not.} \\
\texttt{But I know socks are 0 to me. So the other player *must* want the socks to get anything.} \\
\texttt{If I give the other player the socks, they get points.} \\
\texttt{If I give the other player canvases, they might get points.} \\
\texttt{Let's try to give the other player} \\
            }
        }
    }
    & & \\ \\

    \theutterance \stepcounter{utterance}  
    & & & \multicolumn{2}{p{0.3\linewidth}}{
        \cellcolor[rgb]{0.9,0.9,0.9}{
            \makecell[{{p{\linewidth}}}]{
                \texttt{\tiny{[GM$|$GM]}}
                \texttt{continue} \\
            }
        }
    }
    & & \\ \\

    \theutterance \stepcounter{utterance}  
    & & & \multicolumn{2}{p{0.3\linewidth}}{
        \cellcolor[rgb]{0.9,0.9,0.9}{
            \makecell[{{p{\linewidth}}}]{
                \texttt{\tiny{[GM$|$GM]}}
                \texttt{I will propose a division that maximizes the total value for both of us. Since I value canvases the most (3 points each) and flowers second (2 points each), and socks are worthless to me (0 points), I should aim to get the canvases and flowers for myself, while giving the socks to the other player. This allows the other player to get the socks (which they might value) while I get the high{-}value items.} \\
\\ 
\texttt{Let's assume the other player values socks highly (since they are worthless to me, they are likely the only ones who want them).} \\
\texttt{My plan:} \\
\texttt{{-} I take: 2 flowers, 2 canvases.} \\
\texttt{{-} Other player takes: 3 socks.} \\
\\ 
\texttt{Total value for me: 2*2 + 2*3 = 4 + 6 = 10.} \\
\texttt{Total value for other player: 3 * (value of sock). Even if they value socks at 1, total is 13. If they value them higher, total is even higher.} \\
\texttt{Alternative: I take 2 flowers, 1 canvas. Other takes 3 socks, 1 canvas.} \\
\texttt{My value: 4 + 3 = 7.} \\
\texttt{Other value: 3*sock + 3.} \\
\texttt{Total: 10 + 3 + 3 = 16 (assuming sock=1).} \\
\texttt{My previous plan (I take 2 canvases) gives me 10. If other takes 3 socks, and sock=1, they get 3. Total 13.} \\
\texttt{Wait, if I take 2 canvases and 2 flowers, that's 4 items. The other player gets 3 socks.} \\
\texttt{Is there a better split?} \\
\texttt{If I take 2 canvases (6 pts) and 2 flowers (4 pts) = 10 pts.} \\
\texttt{Other gets 3 socks.} \\
\texttt{If the other player values socks at 1 each, they get 3. Total 13.} \\
\texttt{If I take 1 canvas (3 pts) and 2 flowers (4 pts) = 7 pts.} \\
\texttt{Other gets 2 canvases (6 pts) and 3 socks (3 pts) = 9 pts. Total 16.} \\
\texttt{This seems better if the other player values canvases too. But I don't know their values.} \\
\texttt{However, the rule says "maximize the sum of your score and the score of the other player".} \\
\texttt{If I give the other player the canvases, they might value them highly. If I keep them, they might not.} \\
\texttt{But I know socks are 0 to me. So the other player *must* want the socks to get anything.} \\
\texttt{If I give the other player the socks, they get points.} \\
\texttt{If I give the other player canvases, they might get points.} \\
\texttt{Let's try to give the other player} \\
            }
        }
    }
    & & \\ \\

    \theutterance \stepcounter{utterance}  
    & & \multicolumn{4}{p{0.6\linewidth}}{
        \cellcolor[rgb]{0.9,0.9,0.9}{
            \makecell[{{p{\linewidth}}}]{
                \texttt{\tiny{[GM$\rangle$P2]}}
                \texttt{You are playing a cooperative negotiation game in which you have to agree on how to divide a set of items among you and another player.} \\
\\ 
\texttt{Rules:} \\
\texttt{(a) You and the other player are given a set of items. Each of you is also given a secret value function, representing how much you value each type of object.} \\
\texttt{(b) You exchange messages with the other player to agree on who gets which items. You can send a maximum of 5 messages each, or terminate early by making a secret proposal at any time.} \\
\texttt{(c) You are each asked to submit a secret proposal indicating the items you want formatted in square brackets as follows: "{[}Proposal: <number> <object name>, <number> <object name>, <...>{]}"} \\
\texttt{(d) If your proposals are complementary, i.e., there are enough items to fulfill both proposals, each player is awarded a score based on the sum of values for the items they received. Otherwise, both of you get zero points.} \\
\texttt{(e) Your goal is to maximize the sum of your score and the score of the other player.} \\
\\ 
\texttt{Let us start.} \\
\\ 
\texttt{The set of available items is:} \\
\\ 
\texttt{2 flowers, 3 socks, 2 canvases} \\
\\ 
\texttt{Your secret value function is:} \\
\\ 
\texttt{flower: 1, sock: 2, canvas: 1} \\
\\ 
\texttt{IMPORTANT: Your messages, unless it is the secret proposal, are directly transmitted to the other player, so do not include any response to the rules or text announcing your message. To make a secret proposal, use the indicated format. Do not use square brackets when communicating to the other player or it will be interpreted as your secret proposal.} \\
\\ 
\\ 
\texttt{I will propose a division that maximizes the total value for both of us. Since I value canvases the most (3 points each) and flowers second (2 points each), and socks are worthless to me (0 points), I should aim to get the canvases and flowers for myself, while giving the socks to the other player. This allows the other player to get the socks (which they might value) while I get the high{-}value items.} \\
\\ 
\texttt{Let's assume the other player values socks highly (since they are worthless to me, they are likely the only ones who want them).} \\
\texttt{My plan:} \\
\texttt{{-} I take: 2 flowers, 2 canvases.} \\
\texttt{{-} Other player takes: 3 socks.} \\
\\ 
\texttt{Total value for me: 2*2 + 2*3 = 4 + 6 = 10.} \\
\texttt{Total value for other player: 3 * (value of sock). Even if they value socks at 1, total is 13. If they value them higher, total is even higher.} \\
\texttt{Alternative: I take 2 flowers, 1 canvas. Other takes 3 socks, 1 canvas.} \\
\texttt{My value: 4 + 3 = 7.} \\
\texttt{Other value: 3*sock + 3.} \\
\texttt{Total: 10 + 3 + 3 = 16 (assuming sock=1).} \\
\texttt{My previous plan (I take 2 canvases) gives me 10. If other takes 3 socks, and sock=1, they get 3. Total 13.} \\
\texttt{Wait, if I take 2 canvases and 2 flowers, that's 4 items. The other player gets 3 socks.} \\
\texttt{Is there a better split?} \\
\texttt{If I take 2 canvases (6 pts) and 2 flowers (4 pts) = 10 pts.} \\
\texttt{Other gets 3 socks.} \\
\texttt{If the other player values socks at 1 each, they get 3. Total 13.} \\
\texttt{If I take 1 canvas (3 pts) and 2 flowers (4 pts) = 7 pts.} \\
\texttt{Other gets 2 canvases (6 pts) and 3 socks (3 pts) = 9 pts. Total 16.} \\
\texttt{This seems better if the other player values canvases too. But I don't know their values.} \\
\texttt{However, the rule says "maximize the sum of your score and the score of the other player".} \\
\texttt{If I give the other player the canvases, they might value them highly. If I keep them, they might not.} \\
\texttt{But I know socks are 0 to me. So the other player *must* want the socks to get anything.} \\
\texttt{If I give the other player the socks, they get points.} \\
\texttt{If I give the other player canvases, they might get points.} \\
\texttt{Let's try to give the other player} \\
            }
        }
    }
    & \\ \\

    \theutterance \stepcounter{utterance}  
    & & & \multicolumn{4}{p{0.6\linewidth}}{
        \cellcolor[rgb]{0.9,0.9,0.9}{
            \makecell[{{p{\linewidth}}}]{
                \texttt{\tiny{[GM$\langle$P2]}}
                \texttt{I propose a division that maximizes the total value for both of us. Since I value socks the most (2 points each) and flowers and canvases are worth 1 point each, and I assume the other player likely values socks highly as well (since they are worthless to me), I should aim to secure the socks for myself to ensure a high personal score, while offering the other player the flowers and canvases to ensure they get a positive score, thus maximizing the sum.} \\
\\ 
\texttt{My plan:} \\
\texttt{{-} I take: 3 socks.} \\
\texttt{{-} Other player takes: 2 flowers, 2 canvases.} \\
\\ 
\texttt{Total value for me: 3 * 2 = 6.} \\
\texttt{Total value for other player: 2 * 1 + 2 * 1 = 4.} \\
\texttt{Total combined value: 10.} \\
\\ 
\texttt{This is the maximum possible combined value since the total value of all items is 10 (3 socks * 2 + 2 flowers * 1 + 2 canvases * 1 = 6 + 2 + 2 = 10). By giving the other player everything they value (flowers and canvases) and taking the socks, we ensure the proposal is complementary and the total score is maximized.} \\
\\ 
\texttt{Let's assume the other player values socks highly (since they are worthless to me, they are likely the only ones who want them).} \\
\texttt{My plan:} \\
\texttt{{-} I take: 3 socks.} \\
\texttt{{-} Other player takes: 2 flowers, 2 canvases.} \\
\\ 
\texttt{Total value for me: 3 * 2 = 6.} \\
\texttt{Total value for other player: 2 * 1 + 2 * 1 = 4.} \\
\texttt{Total combined value: 10.} \\
\\ 
\texttt{This is the maximum possible combined value since the total value of all items is 10 (3 socks * 2 + 2 flowers * 1 + 2 canvases * 1 = 6 + 2 + 2 = 10). By giving the other player everything they value (flowers and canvases) and taking the socks, we ensure the proposal is complementary and the total score is maximized.} \\
\\ 
\texttt{I propose: {[}Proposal: 3 socks{]}} \\
            }
        }
    }
     \\ \\

    \theutterance \stepcounter{utterance}  
    & & & \multicolumn{2}{p{0.3\linewidth}}{
        \cellcolor[rgb]{0.9,0.9,0.9}{
            \makecell[{{p{\linewidth}}}]{
                \texttt{\tiny{[GM$|$GM]}}
                \texttt{proposal} \\
            }
        }
    }
    & & \\ \\

    \theutterance \stepcounter{utterance}  
    & & & \multicolumn{2}{p{0.3\linewidth}}{
        \cellcolor[rgb]{0.9,0.9,0.9}{
            \makecell[{{p{\linewidth}}}]{
                \texttt{\tiny{[GM$|$GM]}}
                [0, 3, 0]
            }
        }
    }
    & & \\ \\

    \theutterance \stepcounter{utterance}  
    & & \multicolumn{4}{p{0.6\linewidth}}{
        \cellcolor[rgb]{0.9,0.9,0.9}{
            \makecell[{{p{\linewidth}}}]{
                \texttt{\tiny{[P1$\langle$GM]}}
                \texttt{The other player has submitted a secret proposal. It is now your turn to submit a proposal of your own.} \\
            }
        }
    }
    & \\ \\

    \theutterance \stepcounter{utterance}  
    & \multicolumn{4}{p{0.6\linewidth}}{
        \cellcolor[rgb]{0.9,0.9,0.9}{
            \makecell[{{p{\linewidth}}}]{
                \texttt{\tiny{[P1$\rangle$GM]}}
                \texttt{{[}Proposal: 2 flowers, 2 canvases{]}} \\
            }
        }
    }
    & & \\ \\

    \theutterance \stepcounter{utterance}  
    & & & \multicolumn{2}{p{0.3\linewidth}}{
        \cellcolor[rgb]{0.9,0.9,0.9}{
            \makecell[{{p{\linewidth}}}]{
                \texttt{\tiny{[GM$|$GM]}}
                \texttt{proposal} \\
            }
        }
    }
    & & \\ \\

    \theutterance \stepcounter{utterance}  
    & & & \multicolumn{2}{p{0.3\linewidth}}{
        \cellcolor[rgb]{0.9,0.9,0.9}{
            \makecell[{{p{\linewidth}}}]{
                \texttt{\tiny{[GM$|$GM]}}
                [2, 0, 2]
            }
        }
    }
    & & \\ \\

    \theutterance \stepcounter{utterance}  
    & & & \multicolumn{2}{p{0.3\linewidth}}{
        \cellcolor[rgb]{0.9,0.9,0.9}{
            \makecell[{{p{\linewidth}}}]{
                \texttt{\tiny{[GM$|$GM]}}
                [[2, 0, 2], [0, 3, 0]]
            }
        }
    }
    & & \\ \\

\end{supertabular}
}

\end{document}
